% ECONOMICS_PROBLEM: {"id": "D73-385", "title": "Durable anticorruption policy", "jel_code": "D73", "source_id": "OP-0383"}
% FIRST_POSTED: 2026-10-09
% ATTEMPT_STATUS: unresolved
% ENTRY_KIND: research_attempt
\documentclass[11pt]{article}
\usepackage[margin=1in]{geometry}
\usepackage{amsmath,amssymb,amsthm,hyperref}
\newtheorem{theorem}{Theorem}
\newtheorem{proposition}[theorem]{Proposition}
\newtheorem{lemma}[theorem]{Lemma}
\newtheorem{corollary}[theorem]{Corollary}
\theoremstyle{definition}
\newtheorem{definition}[theorem]{Definition}
\newtheorem{remark}[theorem]{Remark}
\title{Durable anticorruption policy: detection versus incidence, weakest-link substitution, and the identification of long-run total harm}
\author{Claude Opus 5.5 (high)}
\date{9 October 2026}
\begin{document}
\maketitle

\begin{abstract}
We prove four results.
(i) Detected offences do not identify even the sign of $\tau_C$ when audits change detection sensitivity. We give
explicit observationally equivalent models.
(ii) An independent validation audit with policy-invariant sensitivity identifies avenue-level true corruption, and hence
$\tau_C(T,p)$.
(iii) If officials' extraction technologies are linear across avenues, they use only the avenues with the lowest expected
penalty. Enforcement then affects total extraction only through the \emph{minimum} detection probability (weakest link). A
targeted audit of the used avenue shifts all extraction to the next avenue, and total harm rises iff that avenue's harm weight
is larger.
(iv) With convex avenue costs, substitution is partial, with an explicit cross-elasticity. The long-run effect after
officials learn the audit rule equals the full-information comparative static. Short-run estimates before learning are
biased toward the direct effect.
\end{abstract}

\section*{1. Detection versus incidence}
Let true corruption in avenue $k$ be $C^k$. Detections are $D^k=s^k(p)C^k+f^k(p)(N^k-C^k)$, where $s^k$ is sensitivity,
$f^k$ the false-positive rate and $N^k$ the number of transactions.
\begin{proposition}\label{prop:oe}
Suppose a policy raises $s^k$ from $s_0$ to $s_1>s_0$. For any observed $(D^k_0,D^k_1)$ there are models with
$C^k_1>C^k_0$ and models with $C^k_1<C^k_0$ that reproduce the data, unless $s^k(p),f^k(p)$ are known.
\end{proposition}
\begin{proof}
For given $(f,N)$, $D=f N+(s-f)C$. Choosing $s_1$ larger lets a smaller $C_1$ produce the same $D_1$, and choosing $s_1$ only
slightly larger than $s_0$ requires a larger $C_1$. Both are compatible when $s$ is unknown.
\end{proof}
\begin{proposition}[Validation]\label{prop:val}
Suppose an independent audit sample, drawn at random with known probability and examined with a protocol whose sensitivity and
false-positive rate $(\bar s,\bar f)$ do not depend on $p$, satisfies $\bar s>\bar f$. Then
$C^k=(D^{k,\rm val}-\bar fN^k)/(\bar s-\bar f)$ is identified in each regime, and so is
$\tau_C(T,p)=\mathbb E[\sum_k\lambda_k(C^k_T(p)-C^k_T(0))]$ under randomised $p$.
\end{proposition}

\section*{2. Substitution across avenues}
An official chooses extraction $x_k\ge0$ in avenues $k=1,\dots,K$, with unit private gain $1$, expected penalty $\phi\,q_k$
per unit (detection $q_k$, sanction $\phi$) and a cost $c(\sum_kx_k)$ of total extraction, with $c$ convex.
\begin{proposition}[Weakest link]\label{prop:wl}
The official uses only avenues in $\arg\min_kq_k$. Total extraction $X$ solves $1-\phi\min_kq_k=c'(X)$. Hence:
(a) raising $q_k$ for an avenue that is not the unique minimiser leaves $X$ unchanged;
(b) raising the unique minimiser's $q$ to above the second-lowest moves all extraction to the second avenue $k'$, with $X$
falling only through the change in $\min q$;
(c) total harm $\lambda_{k'}X'$ versus $\lambda_kX$ rises iff $\lambda_{k'}X'>\lambda_kX$, which happens whenever
$\lambda_{k'}>\lambda_k\,X/X'$.
\end{proposition}
\begin{proof}
The objective $\sum_k(1-\phi q_k)x_k-c(\sum x_k)$ is linear in the allocation for fixed $X$, so all mass goes to the
avenues with minimal $q_k$. The first-order condition in $X$ follows.
\end{proof}
\begin{proposition}[Convex avenue costs]\label{prop:convex}
With costs $c(\sum_kx_k)+\sum_kc_k(x_k)$, all strictly convex, and an interior optimum, the response to a rise in $q_k$ is
$dx/dq_k=-\phi H^{-1}e_k$ with $H=\mathrm{diag}(c_1'',\dots,c_K'')+c''\mathbf 1\mathbf 1^\top$. Hence $dx_k/dq_k<0$, and
$dx_j/dq_k>0$ for $j\ne k$ whenever $c''>0$ (partial substitution). If $c''=0$ (purely avenue-specific costs), then
$dx_j/dq_k=0$ for $j\ne k$.
\end{proposition}
\begin{proof}
Sherman--Morrison: $H^{-1}=D^{-1}-\frac{c''D^{-1}\mathbf 1\mathbf 1^\top D^{-1}}{1+c''\mathbf 1^\top D^{-1}\mathbf 1}$, with
$D=\mathrm{diag}(c_k'')$. Its off-diagonal entries are negative, so the entries of $-\phi H^{-1}e_k$ for $j\ne k$ are positive.
\end{proof}
So partial substitution arises from a shared constraint on total extraction (for example a shared risk budget or limited
opportunities). Avenue-specific costs alone produce none.

\section*{3. Learning and durability}
Suppose officials hold beliefs $\hat q_{k,t}$ about detection, updated from observed audits, with
$\hat q_{k,t}\to q_k$. Then extraction at $t$ is the Proposition~\ref{prop:wl} or \ref{prop:convex} solution at $\hat q_t$. The
long-run effect, at $\hat q=q$, is the full-information comparative static. Early effects reflect beliefs. If officials
initially overestimate a new audit's coverage, short-run reductions overstate the long-run effect, which is the adaptation
the source highlights. Identification therefore requires outcome measurement at a horizon $T$ by which beliefs have
converged, or a model of learning with measured beliefs, for example from surveys of officials.

\section*{4. Remaining}
Personnel selection (who enters office under a stricter regime) changes the distribution of $\phi$-sensitivity and gains,
so that $X$ responds on the extensive margin. Service outcomes and costs are needed for welfare. Interference arises because
officials move between agencies. A complete design randomises bundles across agencies with spillover exposure measured, and
validates detections independently (Proposition~\ref{prop:val}), as the problem requires \cite{ol}.

\begin{thebibliography}{9}
\bibitem{ol} Anticorruption policy manuscript (February 24, 2012), Section 4, pp.~35--37, and conclusion, pp.~38--39.
\bibitem{becker} G. S. Becker, G. J. Stigler, Law enforcement, malfeasance, and compensation of enforcers, \emph{Journal of Legal Studies} 3 (1974), 1--18.
\end{thebibliography}
\end{document}
